\documentclass[sigconf,nonacm]{acmart}

%% Axiomander -- Iris workshop paper
%% ACM acmart class, sigconf (two-column workshop/conference) format.
%% nonacm: this is a self-archived draft, not a submitted ACM proceedings copy.

\usepackage{booktabs}
\usepackage{mathtools}
\usepackage{mathpartir}
\usepackage{listings}
\usepackage{xcolor}

%% Separation-logic and Iris notation
\newcommand{\sepc}{\mathbin{\ast}}
\newcommand{\pto}{\mapsto}
\newcommand{\wand}{\mathrel{-\!\ast}}
\newcommand{\pure}[1]{\ulcorner #1 \urcorner}
\newcommand{\wpre}{\mathrm{wp}}
\newcommand{\lget}{\mathrm{lget}}
\newcommand{\lupd}{\mathrm{lupd}}
\newcommand{\clobber}{\mathrm{clobber}}
\newcommand{\hoare}[3]{\{#1\}\ #2\ \{#3\}}
\newcommand{\code}[1]{\texttt{#1}}
\let\Bbbk\relax
\usepackage{amssymb}

%% SnakeletLang / SnakeletWp notation
\newcommand{\later}{{\triangleright}}
\newcommand{\wpe}[2]{\mathrm{WPE}\ #1\ \{\!\!\{\, #2 \,\}\!\!\}}
\newcommand{\rval}[1]{\mathsf{val}\,#1}
\newcommand{\rexn}[2]{\mathsf{exn}\,#1\,#2}
\newcommand{\sstep}{\leadsto}
\newcommand{\subs}[3]{#1[#2{:=}#3]}
\newcommand{\entries}{\mathrm{entries}}
\newcommand{\binopd}[2]{\mathrm{binop}(\mathit{op},#1,#2)}

%% Lightweight code listing style for Python / Coq snippets
\lstdefinestyle{snippet}{
  basicstyle=\ttfamily\footnotesize,
  breaklines=true,
  columns=fullflexible,
  keepspaces=true,
  showstringspaces=false,
  frame=single,
  framesep=4pt,
  rulecolor=\color{gray!40},
  xleftmargin=2pt,
  xrightmargin=2pt,
}
\lstset{style=snippet}

\begin{document}

\title{Specification Composition for LLM-Driven Python Verification:
Why We Build on Iris}
\subtitle{A short workshop paper for the Iris community}

\author{Gavin Mendel-Gleason}
\affiliation{%
  \institution{Scidonia}
  \country{}
}
\email{}

\renewcommand{\shortauthors}{Mendel-Gleason}

\begin{abstract}
Axiomander is a verification pipeline that turns ordinary, decorator-free Python
into machine-checked proof obligations and discharges them through a ladder of
mechanical tactics, SMT, and an LLM proof oracle. Its purpose is
\emph{vericoding}: a human iterates on a strong specification, an LLM derives the
implementation (and the helper specifications needed to prove it), and the
verifier delivers a deterministic verdict with a structured residual on failure.
The viability of this loop rests entirely on \emph{specification composition} --
the ability to verify a caller against a callee's \emph{contract} rather than its
body, to reuse proofs across edits, and to stub unverified libraries by their
declared behaviour.

This paper argues that \textbf{separation logic, as realised by Iris, is the
right foundation for that composition}, and reports our experience replacing an
ad-hoc framing mechanism with a native Iris one. Our first backend lowered Python
to a small imperative language (IMP) over a flat store and recovered locality
through generated per-callee, per-variable \emph{frame lemmas} over an explicit
\code{clobber} operator. That mechanism works but scales badly: it produces
WP-term blowup and brittle compiled-Ltac pattern matching, and it does not extend
to aliasing, ownership transfer, or concurrency. We describe a second backend,
\textbf{SnakeletLang/SnakeletWp}, an Iris language and weakest-precondition
calculus (currently 0 \code{Admitted}) in which a callee's footprint is a
separating resource, framing is the separating conjunction, and a \emph{contract
table} lets calls be treated opaquely (as stubs) or transparently (unfolded). We
give Snakelet's full term language and weakest-precondition rule set, the
contract-to-\code{iProp} compilation, the opaque-call rule that
underpins stubbing, and a staged, syntax-directed proof generator whose stages
fail independently and feed residual goals to SMT and an LLM. We close with the
open questions we would most like the Iris community's input on.
\end{abstract}

\keywords{separation logic, Iris, weakest preconditions, specification
composition, stubbing, frame rule, LLM-assisted proof, Python}

\maketitle

\section{The setting: vericoding needs composition}

Axiomander's development model is a loop. The human writes a \emph{strong
specification} -- a precondition/postcondition pair -- as the primary artifact.
An LLM proposes an implementation. The verifier reduces ``does the code meet the
spec?'' to a proof obligation and returns either a certificate or a
\emph{structured residual} (the remaining goal with its hypotheses, plus a
counterexample where one exists). The LLM repairs against the residual, not
against the original natural-language prompt. Crucially, the LLM has freedom in
\emph{sub-specification}: to prove an entry point it may introduce helper
functions, each with its own contract, which are checked the same way.

For this loop to be usable on real codebases, three forms of composition must
hold simultaneously:

\begin{enumerate}
  \item \textbf{Modular verification.} A caller must be provable from a callee's
  \emph{contract} alone. If proving \code{f} required unfolding the bodies of
  everything \code{f} transitively calls, neither proofs nor LLM prompts would
  scale.
  \item \textbf{Proof reuse under iteration.} Editing a function body must not
  invalidate its callers' proofs as long as the function's \emph{contract} is
  unchanged. The verifier should behave like an incremental build system:
  \emph{body changes invalidate local proofs; contract changes invalidate
  callers.}
  \item \textbf{Stubbing.} Libraries we cannot or will not verify must be
  replaceable by their declared behaviour. A caller proves against the stub's
  contract, which becomes an assumption; the trust boundary is the stub, written
  and reviewed once.
\end{enumerate}

All three are, at bottom, the \textbf{frame property}: reasoning about one
component must be insensitive to the parts of the state it does not touch. This is
exactly what separation logic makes primitive, and it is why -- after building
framing the hard way once -- we are migrating Axiomander's proof backend to Iris.

\section{First backend: framing by hand over a flat store}
\label{sec:imp}

Our original backend lowers a verified subset of Python to a small imperative
language IMP, with a state that is a flat finite map from variable names to values
(\code{VZ\,|\,VBool\,|\,VUnit}, plus structural \code{VList\,|\,VTuple\,|\,VDict},
with mutation modelled on a parallel heap representation). The WP calculus is
standard and proven sound in Rocq:
\begin{align*}
\wpre(\mathsf{skip}, Q) &= Q, \\
\wpre(x := e, Q) &= Q[x \mapsto e], \\
\wpre(c_1; c_2, Q) &= \wpre(c_1, \wpre(c_2, Q)), \\
\wpre(\mathsf{if}\ e\ c_1\ c_2, Q) &= (e \Rightarrow \wpre(c_1,Q)) \\
  &\quad {}\wedge (\neg e \Rightarrow \wpre(c_2,Q)).
\end{align*}

Calls are the interesting case. A call site \code{t := g(a)} is verified against
\code{g}'s contract using a \code{CCall} rule. To recover locality -- to let a
caller conclude that variables it cares about survive the call -- we model the
callee's effect with an explicit \code{clobber} operator that havocs the callee's
declared write set, and we must discharge, for \emph{every} variable $v$, an
obligation of the form
\begin{multline*}
\forall v.\; v \notin (\textit{target} :: \textit{writes}) \Rightarrow \\
\lget\,s\,v = \lget\,(\clobber\,(\lupd\,s\,\textit{target}\,r)\,\textit{writes})\,v.
\end{multline*}

This is the frame rule, but encoded extensionally over a flat store. Two problems
follow, both of which we hit in practice:

\begin{itemize}
  \item \textbf{WP-term blowup.} Each successive call nests the full WP expansion
  inside the previous call's postcondition. A function with a handful of
  sequential calls produces terms too large for the kernel to handle comfortably.
  \item \textbf{Brittle automation.} The single universally-quantified frame
  subgoal does not match cleanly in compiled Ltac, partly because coercions
  (e.g.\ an \code{ls} wrapper) normalise away in the \code{.vo} and defeat the
  source-level pattern.
\end{itemize}

Our mitigation was to generate, at the Python IR level, \textbf{one frame lemma
per (callee, preserved-variable) pair}:

\begin{lstlisting}
(* g writes ["result"]; the caller's "a" is
   preserved across t := g(a) *)
Lemma g_frame_a : forall (s : state) (r : Z),
  ~ In "a" ("t" :: "result" :: nil) ->
  lget s "a" =
    lget (clobber (lupd s "t" (VZ r))
                  ("result" :: nil)) "a".
Proof. apply wp_ccall_frame. Qed.
\end{lstlisting}

The caller's proof then replaces the monolithic $\forall v$ subgoal with a
sequence of \code{apply g\_frame\_a. apply g\_frame\_b. ...}, each trivial and
each failing independently with its own residual. This restores tractability and
keeps the per-obligation, fail-locally discipline we want. But it is, frankly, a
workaround: we are \emph{reconstructing} the frame rule, one variable at a time,
in a logic that does not have it. It does not model aliasing, it does not express
ownership transfer, and it has no path to concurrency. The lesson is the familiar
one: \textbf{if framing is load-bearing, the assertion logic should be
substructural.}

\section{Second backend: SnakeletLang and SnakeletWp over Iris}
\label{sec:iris}

We have built a second backend in which framing is not reconstructed but
inherited. \textbf{Snakelet} is a small Iris \emph{intermediate language} for the
verified Python subset, and \textbf{SnakeletWp} is its weakest-precondition
calculus, defined inside the Iris program logic. The state is an Iris heap; a
mutable Python object field is a points-to assertion $l \pto v$, and ownership is
a separating resource rather than an entry in a flat map; assertions are
\code{iProp}. The next three subsections give Snakelet's term language
(Fig.~\ref{fig:grammar}), its operational semantics (Fig.~\ref{fig:opsem}), and
its weakest-precondition rules (Fig.~\ref{fig:wp}); the calculus currently carries
\textbf{0 \code{Admitted}} across the entire rule set -- values, binary operators,
\texttt{let}, \texttt{if}, load/store/alloc, the two call forms, and the exception
and loop rules whose postcondition ranges over a \code{Result}.

\subsection{Snakelet: an intermediate language for the verified subset}
\label{sec:snakelet-il}

Axiomander does not reason about arbitrary Python. The front end lowers a
\emph{verified subset} -- straight-line assignment, conditionals, \code{while} and
\code{for} loops, \code{try}/\code{except}, and calls -- to Snakelet, and every
proof obligation is discharged at the Snakelet level. The Python-to-Snakelet
lowering is the trusted boundary of the pipeline (intended to be extracted from
Coq later); Snakelet itself is small enough to carry a complete operational
semantics and WP calculus, which is what this section gives.

Snakelet is an expression language in administrative-normal form: every operator
and call takes values or variables, and evaluation order is made explicit by
\code{let}-binding and a per-item evaluation-context grammar. A store is a finite
map from locations to values ($\sigma : \mathrm{loc} \rightharpoonup v$); a mutable
Python object field is a heap cell, exposed in the logic as $l \pto v$. The
immutable Python containers -- list, tuple, dict, set -- are \emph{structural}
values, not heap objects, mirroring Dafny's \code{seq} and F\textsuperscript{$\star$}'s
\code{list}.

\paragraph{Two coordinated developments.} The core language carries float
literals, dictionary get/set, and the concurrency primitives \code{FAA}/\code{fork};
the exception-aware extension presented here adds first-class exceptions and
carries the complete, machine-checked rule set. We present the latter because
exceptions change the \emph{shape of the judgement}. Following van
Collem/de~Vilhena/Krebbers (PLDI~2026), $\mathsf{raise}\;v$ on a value is a
\textbf{stuck terminal} expression -- not itself a value -- so an uncaught raise
sits in evaluation position as the program's exceptional result. A program thus
terminates in one of two ways, captured by the result type $r ::= \mathsf{val}\;v
\mid \mathsf{exn}\;s\;v$, and (Section~\ref{sec:snakelet-wp}) the WP postcondition
ranges over $r$ rather than over values. Figure~\ref{fig:grammar} gives the term
language.

\begin{figure}[t]
\footnotesize
\[
\begin{array}{r@{\ }c@{\ }l}
v &::=& n \mid b \mid \mathit{fl} \mid s \mid l \mid () \mid \mathsf{exn}\;s\;v \\
  &\mid& [\,\overline{v}\,] \mid (\,\overline{v}\,) \mid \{\,\overline{v\!:\!v}\,\} \mid \{\,\overline{v}\,\} \\[2pt]
\mathit{op} &::=& {+} \mid {-} \mid {\times} \mid {/} \mid \mathbin{\%} \mid {=} \mid {\neq} \mid {<} \mid {\leq} \mid {>} \mid {\geq} \\
  &\mid& \mathsf{and} \mid \mathsf{or} \mid \mathsf{in} \mid \mathsf{len} \mid \mathsf{union} \mid \mathsf{inter} \\[2pt]
e &::=& v \mid x \mid \mathsf{let}\;x{=}e\;\mathsf{in}\;e \mid \mathit{op}(e,e) \mid g(\overline{e}) \\
  &\mid& \mathord{!}e \mid e \mathbin{:=} e \mid \mathsf{ref}\;e \mid \mathsf{if}\;e\;\mathsf{then}\;e\;\mathsf{else}\;e \\
  &\mid& \mathsf{raise}\;e \mid \mathsf{try}\;e\;\mathsf{catch}\;x{\Rightarrow}e \\
  &\mid& \mathsf{while}\;e\;e \mid \mathsf{for}\;x\;\mathsf{in}\;e\;\mathsf{do}\;e \\[2pt]
K &::=& \mathsf{let}\;x{=}\square\;\mathsf{in}\;e \mid \mathit{op}(\square,v) \mid \mathit{op}(e,\square) \mid \mathord{!}\square \mid \mathsf{ref}\;\square \\
  &\mid& \square \mathbin{:=} v \mid e \mathbin{:=} \square \mid \mathsf{if}\;\square\;\mathsf{then}\;e\;\mathsf{else}\;e \mid \mathsf{raise}\;\square \\
  &\mid& \mathsf{try}\;\square\;\mathsf{catch}\;x{\Rightarrow}e \mid \mathsf{for}\;x\;\mathsf{in}\;\square\;\mathsf{do}\;e \\[2pt]
r &::=& \mathsf{val}\;v \mid \mathsf{exn}\;s\;v
\end{array}
\]
\caption{Snakelet term language: values $v$ (integer, boolean, float, string,
location, unit, exception object, and the structural list/tuple/dict/set),
operators $\mathit{op}$, expressions $e$, single-item evaluation contexts $K$, and
results $r$. Every $K$ is \emph{neutral} except $\mathsf{try}\;\square\,\dots$;
$\mathit{op}(e,e)$ evaluates its right operand first. The core development
additionally provides $\mathsf{FAA}$, $\mathsf{fork}$, and dictionary get/set.}
\label{fig:grammar}
\end{figure}

\subsection{Operational semantics}
\label{sec:snakelet-op}

Reduction is small-step and factors into \emph{pure} steps $e \sstep e'$ (no heap,
deterministic) and \emph{head} steps $(e,\sigma) \rightarrow (e',\sigma';
\vec{\mathit{ef}})$ (heap and calls), each lifted under an evaluation context by
the two decomposition rules \textsc{Ctx-P}/\textsc{Ctx-H} (Figure~\ref{fig:opsem}).
Contexts are single-item because the WP binds one frame at a time.

All context items are \textbf{neutral} except $\mathsf{try}\;\square\;\mathsf{catch}
\;x{\Rightarrow}h$. The unwinding rule \textsc{Unwind}, $K[\mathsf{raise}\;v]
\sstep \mathsf{raise}\;v$, fires only for neutral $K$, so a raise propagates
through \code{let}, operators, loads, and the rest, but is \emph{caught} by an
enclosing \code{try} frame (\textsc{TryC}) -- exactly Python/ML exception
propagation. Calls dispatch through a \emph{contract table} $\entries$: a name maps
to either $\mathsf{FunSpec}\;P\;Q$ (\emph{opaque} -- the call steps to some result
satisfying $Q$ only when $P$ holds, so calling outside the precondition is stuck)
or $\mathsf{FunDef}\;\vec{x}\;b$ (\emph{transparent} -- substitute the arguments
into the body). The two modes share one table with one entry per name, so they are
mutually exclusive by construction.

\begin{figure*}[t]
\centering
\small
\textbf{Pure reduction}\quad $e \sstep e'$
\begin{mathpar}
\inferrule*[left=\textsc{Let}]{\,}{\mathsf{let}\;x{=}v\;\mathsf{in}\;e \sstep \subs{e}{x}{v}}
\and
\inferrule*[left=\textsc{Op}]{\,}{\mathit{op}(v_1,v_2) \sstep \binopd{v_1}{v_2}}
\and
\inferrule*[left=\textsc{IfT}]{\,}{\mathsf{if}\;\mathsf{true}\;\mathsf{then}\;e_1\;\mathsf{else}\;e_2 \sstep e_1}
\and
\inferrule*[left=\textsc{IfF}]{\,}{\mathsf{if}\;\mathsf{false}\;\mathsf{then}\;e_1\;\mathsf{else}\;e_2 \sstep e_2}
\and
\inferrule*[left=\textsc{TryV}]{\,}{\mathsf{try}\;v\;\mathsf{catch}\;x{\Rightarrow}h \sstep v}
\and
\inferrule*[left=\textsc{TryC}]{\,}{\mathsf{try}\;(\mathsf{raise}\;v)\;\mathsf{catch}\;x{\Rightarrow}h \sstep \subs{h}{x}{v}}
\and
\inferrule*[left=\textsc{Unwind}]{\mathrm{neutral}(K)}{K[\mathsf{raise}\;v] \sstep \mathsf{raise}\;v}
\and
\inferrule*[left=\textsc{While}]{\,}{\mathsf{while}\;e_1\;e_2 \sstep \mathsf{if}\;e_1\;\mathsf{then}\;(e_2 \mathbin{;} \mathsf{while}\;e_1\;e_2)\;\mathsf{else}\;()}
\and
\inferrule*[left=\textsc{ForNil}]{\,}{\mathsf{for}\;x\;\mathsf{in}\;[\,]\;\mathsf{do}\;b \sstep ()}
\and
\inferrule*[left=\textsc{ForCons}]{\,}{\mathsf{for}\;x\;\mathsf{in}\;(v\!::\!\mathit{vs})\;\mathsf{do}\;b \sstep \subs{b}{x}{v} \mathbin{;} \mathsf{for}\;x\;\mathsf{in}\;\mathit{vs}\;\mathsf{do}\;b}
\end{mathpar}
\smallskip
\textbf{Head reduction}\quad $(e,\sigma) \rightarrow (e',\sigma';\vec{\mathit{ef}})$ \qquad and \quad \textbf{context decomposition}
\begin{mathpar}
\inferrule*[left=\textsc{Load}]{\sigma(l)=v}{(\mathord{!}l,\sigma) \rightarrow (v,\sigma)}
\and
\inferrule*[left=\textsc{Store}]{l \in \mathrm{dom}\,\sigma}{(l \mathbin{:=} v,\sigma) \rightarrow ((),\,\sigma[l{\mapsto}v])}
\and
\inferrule*[left=\textsc{Alloc}]{\sigma(l)=\bot}{(\mathsf{ref}\;v,\sigma) \rightarrow (l,\,\sigma[l{\mapsto}v])}
\and
\inferrule*[left=\textsc{CallSpec}]{\entries(g)=\mathsf{FunSpec}\;P\;Q \\ P\,\vec{v} \\ Q\,\vec{v}\,v}{(g(\vec{v}),\sigma) \rightarrow (v,\sigma)}
\and
\inferrule*[left=\textsc{CallDef}]{\entries(g)=\mathsf{FunDef}\;\vec{x}\;b \\ |\vec{v}|=|\vec{x}|}{(g(\vec{v}),\sigma) \rightarrow (\subs{b}{\vec{x}}{\vec{v}},\sigma)}
\and
\inferrule*[left=\textsc{Ctx-P}]{e \sstep e'}{(K[e],\sigma) \rightarrow (K[e'],\sigma)}
\and
\inferrule*[left=\textsc{Ctx-H}]{(e,\sigma) \rightarrow (e',\sigma';\vec{\mathit{ef}})}{(K[e],\sigma) \rightarrow (K[e'],\sigma';\vec{\mathit{ef}})}
\end{mathpar}
\caption{Snakelet operational semantics. Pure steps are heap-independent and
deterministic; head steps touch the heap or the contract table $\entries$.
\textsc{CallSpec} fires only inside the callee's precondition $P$ (calling outside
$P$ is stuck), making the precondition a proof obligation on the caller.
$\vec{\mathit{ef}}$ collects forked threads (empty except for \code{fork} in the
core development).}
\label{fig:opsem}
\end{figure*}

\subsection{The weakest-precondition rules}
\label{sec:snakelet-wp}

SnakeletWp is a hand-rolled weakest precondition whose postcondition is $\Phi :
\code{Result} \to \mathit{iProp}$. It is the guarded fixpoint of: if $e$ is
terminal ($\mathrm{result}(e) = \mathsf{Some}\,r$) then $\Phi\,r$ holds under a
fancy update; otherwise $e$ is reducible in the current heap and, after one step,
the state interpretation (the Iris \code{gen\_heap} authoritative view) and the WP
of the reduct are re-established. Ranging $\Phi$ over \code{Result} is the one
structural change exceptions force; everything else is standard Iris.
Figure~\ref{fig:wp} lists the rules, all proved with \textbf{0 \code{Admitted}}.

Three points. First, \textbf{composition goes through} $\mathrm{bindp}$: the bind
rule \textsc{Bind} threads a \code{Result}-typed continuation
$\mathrm{bindp}\,K\,\Phi$ defined by $\mathrm{bindp}\,K\,\Phi\,(\rval{v}) =
\wpe{K[v]}{\Phi}$ and $\mathrm{bindp}\,K\,\Phi\,(\rexn{s}{v}) = \Phi(\rexn{s}{v})$
-- on a value it keeps evaluating the context, on an exception it short-circuits to
the exceptional postcondition. This is the convergence-critical lemma: it is what
makes the exception WP \emph{compose}, and with \textsc{Unwind}/\textsc{LetRaise}
it delivers full propagation. Second, the \textbf{heap rules are orthogonal to
exceptions}: \textsc{Load}/\textsc{Store}/\textsc{Alloc} are the usual Iris
points-to rules, and the exceptional arm of $\Phi$ describes the heap \emph{as it
stood at the raise} -- a function that mutates a cell then raises satisfies a
postcondition whose $\rexn{s}{}$ arm asserts $l \pto v$ with the \emph{mutated}
value. Third, the \textbf{two call rules are the composition story} expanded in
Section~\ref{sec:opaque}: \textsc{CallOpaque} is stubbing and modular verification
(the caller establishes $P\,\vec{v}$, receives $Q\,\vec{v}\,v$, and never sees the
body); \textsc{CallDef} unfolds a first-party body.

For loops, \textsc{While}/\textsc{ForNil}/\textsc{ForCons} unfold one iteration. On
top of \textsc{ForCons} we derive a structural fold rule that carries a user
invariant $P$ over the \emph{remaining suffix} of the iterated list and
re-establishes it (or escapes through the exceptional arm) at each element:
\[
\begin{array}{l}
P\,M \;\sepc\; \big(\,\Box\,\forall v\,\mathit{vs}.\ P\,(v\!::\!\mathit{vs}) \wand
  \wpe{\subs{b}{x}{v}}{r.\,\theta(r)}\,\big) \\[2pt]
\quad {}\sepc\; (P\,[\,] \wand \Phi(\rval{()})) \;\;\vdash\;\;
  \wpe{\mathsf{for}\;x\;\mathsf{in}\;M\;\mathsf{do}\;b}{\Phi},
\end{array}
\]
where $\theta(\rval{\_}) = P\,\mathit{vs}$ and $\theta(\rexn{s}{v}) =
\Phi(\rexn{s}{v})$; specialising $P\,M := \pure{\mathrm{Forall}\,Q\,M}$ yields the
per-element \code{forall} loop the generator emits. A dedicated heap-counter
\code{while} rule closes counting loops by L\"ob induction. Each rule has a
matching \emph{stage tactic} (Section~\ref{sec:staged}): \code{pure\_step} fires
\textsc{Let}\slash\textsc{Op}\slash\textsc{IfT}\slash\textsc{IfF}\slash\textsc{TryV}\slash\textsc{TryC},
\code{heap\_load}\slash\code{store}\slash\code{alloc} the heap rules,
\code{call\_opaque}\slash\code{call\_transparent} the call rules, \code{raise\_step}
\textsc{Raise}\slash\textsc{Unwind}, and \code{loop\_unfold} the loop rules.

\begin{figure*}[t]
\centering
\small
\begin{mathpar}
\inferrule*[left=\textsc{Val}]{\,}{\Phi(\rval{v}) \vdash \wpe{v}{\Phi}}
\and
\inferrule*[left=\textsc{Raise}]{\,}{\Phi(\rexn{s}{v}) \vdash \wpe{\mathsf{raise}\;(\rexn{s}{v})}{\Phi}}
\and
\inferrule*[left=\textsc{Wand}]{\,}{\wpe{e}{\Phi} \sepc (\forall r.\ \Phi\,r \wand \Psi\,r) \vdash \wpe{e}{\Psi}}
\and
\inferrule*[left=\textsc{Let}]{\,}{\later \wpe{\subs{e}{x}{v}}{\Phi} \vdash \wpe{\mathsf{let}\;x{=}v\;\mathsf{in}\;e}{\Phi}}
\and
\inferrule*[left=\textsc{LetRaise}]{\,}{\Phi(\rexn{s}{v}) \vdash \wpe{\mathsf{let}\;x{=}\mathsf{raise}\,(\rexn{s}{v})\;\mathsf{in}\;e}{\Phi}}
\and
\inferrule*[left=\textsc{Op}]{\,}{\later \wpe{\binopd{v_1}{v_2}}{\Phi} \vdash \wpe{\mathit{op}(v_1,v_2)}{\Phi}}
\and
\inferrule*[left=\textsc{IfT}]{\,}{\later \wpe{e_1}{\Phi} \vdash \wpe{\mathsf{if}\;\mathsf{true}\;\mathsf{then}\;e_1\;\mathsf{else}\;e_2}{\Phi}}
\and
\inferrule*[left=\textsc{IfF}]{\,}{\later \wpe{e_2}{\Phi} \vdash \wpe{\mathsf{if}\;\mathsf{false}\;\mathsf{then}\;e_1\;\mathsf{else}\;e_2}{\Phi}}
\and
\inferrule*[left=\textsc{TryV}]{\,}{\later \wpe{v}{\Phi} \vdash \wpe{\mathsf{try}\;v\;\mathsf{catch}\;x{\Rightarrow}h}{\Phi}}
\and
\inferrule*[left=\textsc{TryC}]{\,}{\later \wpe{\subs{h}{x}{v}}{\Phi} \vdash \wpe{\mathsf{try}\;(\mathsf{raise}\;v)\;\mathsf{catch}\;x{\Rightarrow}h}{\Phi}}
\and
\inferrule*[left=\textsc{Bind}]{\mathrm{neutral}(K)}{\wpe{e}{\mathrm{bindp}\,K\,\Phi} \vdash \wpe{K[e]}{\Phi}}
\and
\inferrule*[left=\textsc{Load}]{\,}{l \pto v \;\sepc\; \later (l \pto v \wand \Phi(\rval{v})) \vdash \wpe{\mathord{!}l}{\Phi}}
\and
\inferrule*[left=\textsc{Store}]{\,}{l \pto w \;\sepc\; \later (l \pto v \wand \Phi(\rval{()})) \vdash \wpe{l \mathbin{:=} v}{\Phi}}
\and
\inferrule*[left=\textsc{Alloc}]{\,}{\later (\forall l.\ l \pto v \wand \Phi(\rval{l})) \vdash \wpe{\mathsf{ref}\;v}{\Phi}}
\and
\inferrule*[left=\textsc{ForNil}]{\,}{\later \Phi(\rval{()}) \vdash \wpe{\mathsf{for}\;x\;\mathsf{in}\;[\,]\;\mathsf{do}\;b}{\Phi}}
\and
\inferrule*[left=\textsc{CallOpaque}]{\entries(g)=\mathsf{FunSpec}\;P\;Q \\ P\,\vec{v}}{\later (\forall v.\ \pure{Q\,\vec{v}\,v} \wand \Phi(\rval{v})) \vdash \wpe{g(\vec{v})}{\Phi}}
\and
\inferrule*[left=\textsc{CallDef}]{\entries(g)=\mathsf{FunDef}\;\vec{x}\;b \\ |\vec{v}|=|\vec{x}|}{\later \wpe{\subs{b}{\vec{x}}{\vec{v}}}{\Phi} \vdash \wpe{g(\vec{v})}{\Phi}}
\end{mathpar}
\[
\begin{array}{l}
\textsc{While}\quad \later \wpe{\mathsf{if}\;e_1\;\mathsf{then}\;(e_2 \mathbin{;} \mathsf{while}\;e_1\;e_2)\;\mathsf{else}\;()}{\Phi} \\
\qquad\quad {}\vdash\; \wpe{\mathsf{while}\;e_1\;e_2}{\Phi} \\[4pt]
\textsc{ForCons}\quad \later \wpe{\subs{b}{x}{v} \mathbin{;} \mathsf{for}\;x\;\mathsf{in}\;\mathit{vs}\;\mathsf{do}\;b}{\Phi} \\
\qquad\quad {}\vdash\; \wpe{\mathsf{for}\;x\;\mathsf{in}\;(v\!::\!\mathit{vs})\;\mathsf{do}\;b}{\Phi}
\end{array}
\]
\caption{SnakeletWp weakest-precondition rules. The postcondition $\Phi :
\code{Result} \to \mathit{iProp}$ ranges over $r ::= \rval{v} \mid \rexn{s}{v}$;
$\later$ is the Iris later modality, $\sepc$ the separating conjunction, $\wand$
the magic wand, and $\pure{\cdot}$ the persistent embedding of a pure proposition.
\textsc{Bind} uses the \code{Result}-transformer $\mathrm{bindp}$ (text); the
wider \textsc{While}/\textsc{ForCons} rules are shown beneath the grid. Every
rule is machine-checked.}
\label{fig:wp}
\end{figure*}

\subsection{Contracts compile to \code{iProp}, composed with the separating
conjunction}

Contracts remain ordinary Python -- plain \code{assert} statements (leading $=$
precondition, trailing-before-return $=$ postcondition, in-loop $=$ invariant) or
a verifier-only \code{axiomander:} docstring. A shared \code{ContractLinter}
parses them into a contract IR; a backend-specific pass (\code{contract\_ir\_iris})
compiles that IR to \code{iProp}. Pure facts are injected with the persistent
embedding $\pure{P}$; resource facts compile to points-to assertions.

\begin{table}[t]
\caption{Contract fragments compiled to \code{iProp}.}
\label{tab:compile}
\small
\begin{tabular}{@{}ll@{}}
\toprule
Python contract fragment & \code{iProp} \\
\midrule
\code{assert x >= 1} & $\pure{x \ge 1}$ \\
\code{assert implies(A, B)} & $\pure{A \to B}$ \\
\code{assert all(x != i for i in range(0,10))} &
  $\pure{\forall i,\ 0 \le i < 10 \to x \ne i}$ \\
\code{owns(box)} & $\textit{box.value} \pto v$ (a resource) \\
\bottomrule
\end{tabular}
\end{table}

The decisive design choice is that \textbf{the default composition of contract
clauses is the separating conjunction $\sepc$, not $\wedge$.} When a precondition
or postcondition is built from several \code{assert}s, they are joined with
$\sepc$. For the current pure-integer fragment this costs nothing, since
$\pure{P} \sepc \pure{Q} \dashv\vdash \pure{P \wedge Q}$; but it means the
\emph{same} composition rule already accounts for disjoint resources, and it is
the connective that will let us scale to invariants and atomic updates without
changing the front end.

A typical mixed pure/resource contract, written by the user with no Iris
notation:

\begin{lstlisting}[language=Python]
def bump(box: Box) -> int:
    """
    axiomander:
        requires:
            owns(box)
            box.value >= 0
        modifies:
            box.value
        ensures:
            owns(box)
            box.value == old(box.value) + 1
            result == box.value
    """
    box.value = box.value + 1
    return box.value
\end{lstlisting}

compiles to the Hoare triple
\begin{multline*}
\{\,\textit{box.value} \pto n \sepc \pure{n \ge 0}\,\}\quad
\textit{bump(box)} \\
\{\,r.\ \textit{box.value} \pto n{+}1 \sepc \pure{r = n{+}1}\,\}.
\end{multline*}

\code{owns(box)} plus \code{modifies: box.value} \emph{is} the resource
footprint; the user never writes $\sepc$, $\pto$, masks, or proof-mode tactics.

\subsection{The frame rule is free, and so is composition}

Because assertions are \code{iProp}, the structural frame rule
\[
\inferrule{\hoare{P}{c}{Q}}{\hoare{P \sepc R}{c}{Q \sepc R}}
\]
holds without a side condition: any resource $R$ describing a disjoint region of
the heap is carried across $c$ by the separating conjunction, automatically.
Compared with Section~\ref{sec:imp}, the entire \code{clobber}/$\forall v \notin
\textit{writes}$ apparatus, and the per-callee per-variable frame-lemma generator
that papered over it, simply disappear. Locality stops being something we encode
and becomes the shape of the logic. This is the single biggest reason for the
migration: \emph{composability is sound, not merely convenient.}

\subsection{Opaque vs.\ transparent calls: stubbing as a logical primitive}
\label{sec:opaque}

Calls are dispatched through a \textbf{contract table} (an Iris \code{FunCtx})
that maps each callee to its specification. SnakeletWp provides two call rules:

\begin{itemize}
  \item \textbf{Transparent call} -- the callee's body is available and is
  unfolded; used for first-party functions being verified in the same run.
  \item \textbf{Opaque call} -- only the callee's \emph{contract} is available;
  the caller consumes the callee's precondition (giving up the relevant
  resources), receives its postcondition, and frames the rest. Schematically, for
  a callee with spec $\hoare{P_g}{g}{Q_g}$:
\end{itemize}
\[
\inferrule{P_g \sepc R \\\\
  \forall r.\ (Q_g \sepc R) \vdash \wpre\ K\ \{\Phi\}}
{\wpre\ (\mathsf{call}\ g;\ K)\ \{\Phi\}}
\]

The opaque rule is exactly \textbf{stubbing as a logical primitive}. A library
function is given a contract -- in a \code{.pyi} stub or an \code{axiomander:}
docstring -- and that contract enters the \code{FunCtx}; callers prove against it
via the opaque rule, and the body is never required. It is also exactly
\textbf{modular verification under iteration}: a transparent callee can be
re-declared opaque once its contract is stable, and its callers' proofs depend
only on the contract, so a body edit does not disturb them. The \code{FunCtx}
entry's hash is the summary that our incremental cache keys on.

\subsection{Staged, syntax-directed proofs with independent failure}
\label{sec:staged}

SnakeletWp ships a small instruction set of \textbf{stage tactics} --
\code{pure\_step}, \code{case\_bool}, \code{call\_opaque},
\code{call\_transparent}, \code{finish\_pure}, and a precondition-discharge stage
\code{call\_opaque\_pre} -- and the proof generator (\code{iris\_proof\_gen})
emits exactly one stage per IR node, chosen by the node's syntax and the callee's
table entry. The generator performs no symbolic execution and needs no knowledge
of intermediate values: each stage tactic extracts what it needs from the Iris
goal at proof time. Loops are handled by an \code{iL\"ob}-based invariant stage,
with the invariant taken verbatim from the user's in-loop \code{assert}. The
proof script \emph{is} the trace.

The payoff for the LLM loop is that \textbf{stages fail independently with
classifiable errors}. When \code{call\_opaque\_pre} cannot discharge a callee
precondition (e.g.\ nonlinear arithmetic, a string fact), the residual is exported
to SMT; an UNSAT result is imported as \code{Axiom smt\_ax\_N : ...} and the stage
is regenerated as \code{call\_opaque\_pre (exact smt\_ax\_N)}. A pure
postcondition that \code{finish\_pure} cannot close (\code{reflexivity}/\code{lia})
is likewise an SMT candidate, and beyond SMT the \emph{structured residual goal}
-- a named stage plus the goal state from coq-lsp, rather than a stderr regex --
is what the LLM oracle receives. Small named obligations give better SMT
performance, better LLM prompts, per-stage caching, and parallelism, and they mean
a failed attempt yields a reusable artifact rather than a dead end.

\subsection{What the Iris choice does and does not change}

The migration is deliberately confined to the proof backend. The
\code{ContractLinter}, the SMT export (\code{smt\_export}, \code{theory\_smt}),
the LLM oracle, and the function-level incremental cache are all backend-agnostic;
the contract IR diverges only at the final rendering step (flat-store
\code{to\_coq} vs.\ \code{contract\_ir\_iris.iris\_prop}). Iris \emph{adds}
per-stage hashing for finer-grained reuse but does not replace the cache. From the
user's side, nothing changes: contracts are still decorator-free Python, and
\code{owns(...)} is the only new, entirely optional, resource vocabulary.

\section{Discussion and relation to prior work}

Axiomander's separation-logic interpretation of object fields and its opaque call
rule are textbook Iris; the novelty is not in the logic but in the \emph{system
context}. To our knowledge the combination is unusual: (i) the assertion logic is
never exposed to the user -- contracts are plain Python and the separating
conjunction is the implicit default composition; (ii) the proof obligations are
generated and largely discharged by an LLM operating on structured residual goals,
with Iris/Rocq as the trust kernel; and (iii) the same \code{FunCtx} opaque rule
serves simultaneously as modular verification, incremental-reuse boundary, and
library stubbing mechanism. We see Axiomander as a stress test of the claim that
Iris's framing and ownership story is the right substrate for \emph{automated},
composition-heavy verification, not only for expert-driven interactive proof.

Our experience also offers a small, concrete data point: building a verifier
without a substructural logic and then retrofitting locality
(Section~\ref{sec:imp}) is costly and bounded; starting from Iris
(Section~\ref{sec:iris}) makes the very properties our LLM loop depends on --
frame, stub, reuse -- definitional.

\section{Status and open questions for the community}

\paragraph{Status (June 2026).} Snakelet and SnakeletWp are complete and free of
\code{Admitted} across the rule set of Figures~\ref{fig:opsem}--\ref{fig:wp}: the
core fragment (binop, let, if, load/store/alloc, opaque/transparent call) and the
exception-aware extension with its \code{Result}-typed postcondition and loop
rules, whose eight-lemma convergence gate is \code{Qed}. The staged generator, the
end-to-end Python-to-Iris pipeline (sharing the \code{ContractLinter}), the
contract-to-\code{iProp} compiler, and the SMT-axiom escalation slot are in place,
with tens of Iris-specific tests passing alongside the mature flat-store backend.
Remaining work wires the exception calculus into the Python lowering and extends
the fragment to loop invariants over shared disjoint memory, in-place
data-structure mutation, and a typed subset
(strings\slash floats\slash\code{isinstance}\slash\code{None}).

We would value the community's view on:

\begin{enumerate}
  \item \textbf{Ownership inference from \code{modifies}.} We currently derive a
  footprint from explicit \code{owns(...)} plus \code{modifies:}. How far can
  footprints be \emph{inferred} for an LLM-generated body without burdening the
  user, and what is the right failure mode when inference is ambiguous?
  \item \textbf{Specs as the LLM's search space.} The LLM proposes both code and
  helper contracts. Are there idioms (magic wands, representation predicates,
  abstract predicates with \code{FunCtx}-level abstraction) that are especially
  amenable to automatic \emph{generation}, as opposed to automatic \emph{use}?
  \item \textbf{Residual goals as oracle inputs.} We feed a failed stage's goal
  state to an LLM. Which Iris proof-mode presentations of a residual (hypothesis
  naming, spatial/persistent split, masks elided) make the best prompts?
  \item \textbf{Concurrency on the same front end.} Our composition default is
  already $\sepc$. What is the minimal extension of the decorator-free contract
  surface that could express invariants/atomic updates without exposing Iris
  syntax to the user?
  \item \textbf{Trust surface.} With SMT-discharged facts imported as axioms and
  LLM-generated proofs re-checked by Rocq, the kernel remains the authority but
  the axiom slots are a trust boundary. What discipline (e.g.\ logged query
  hashes, replay) would the community consider acceptable for reproducibility?
\end{enumerate}

\section*{Artifact pointers}

This paper summarises material developed in the Axiomander repository: the Iris
migration plan and current state; the Iris backend prototype (the
\code{bump}/\code{owns} example and resource classification); the per-callee frame
lemmas in the flat-store backend (Section~\ref{sec:imp}); the staged proof
engineering notes; the incremental-verification cache design; and the companion
whitepaper and slides motivating vericoding.

\end{document}
